\section{The retired household}\label{sec:model}

A household lives $J$ periods, works until age $J_r$ and is retired after it. Period utility is
CRRA with coefficient $\mu$ and discount factor $\beta$, the gross return is $R=1+r$, and assets
may not fall below a borrowing limit, taken to be zero. In retirement labour earnings are replaced
by a pension.

\paragraph{A lump-sum pension.} If the pension is a constant $p$, the retired household faces the
same problem as the working one with earnings replaced by $p$, so the retired phase is an
income-fluctuation problem in its own right (\lean{withPension}) and the whole of the standard
theory applies to it. Two facts are worth recording. A retiree's human wealth does not depend on
the earnings shock, because the pension does not, and it never exceeds the perpetuity value of the
pension, $H_k\le p/r$ (\lean{stageHumanWealth\_withPension\_const},
\lean{stageHumanWealth\_withPension\_le}). And cash on hand is strictly increasing in assets
(\lean{withPension\_resources\_strictMono}). Both fail once the pension is means-tested.

\paragraph{Three schedules.} Write the pension as a function $p(\cdot)$ of assets and cash on hand
as
\[
  m(a)=Ra+p(a)
\]
(\lean{cash}). We consider three schedules (\lean{lumpSum}, \lean{tapered}, \lean{cliff}):
\begin{itemize}[leftmargin=1.5em]
\item a \emph{lump sum}, $p(a)=p$;
\item a \emph{taper}, $p(a)=\max\{0,\;p-\tau\max\{0,a-\bar a\}\}$, which withdraws the pension at
  rate $\tau$ per unit of assets above a free area $\bar a$ and never pays a negative amount;
\item a \emph{cliff}, $p(a)=p$ for $a\le\bar a$ and $0$ above, the limit of a taper as
  $\tau\to\infty$.
\end{itemize}

\paragraph{Where the test is assessed.} It matters whether the pension paid this period depends on
the assets held this period or on the assets carried into the next. A test on assets \emph{carried
forward} enters only the continuation value, leaving today's budget alone. A test on assets
\emph{currently held} enters today's cash on hand directly. Real asset tests are of the second
kind, assessed each period on current holdings, and the two have different consequences, as
Sections~\ref{sec:cash} and~\ref{sec:assets} show.

\paragraph{One period of the problem.} Throughout we write the household's choice as
\begin{equation}\label{eq:step}
  \max_{x\in[0,m)}\ u(m-x)+W(x),
\end{equation}
with $m$ cash on hand, $x$ next period's assets and $W$ the continuation value, and call $x$
\emph{optimal} at $m$ when it attains the maximum (\lean{IsOptimalSaving}). The feasible set
excludes $x=m$, so consumption is positive. Nothing below assumes anything about $W$ beyond what
is stated: in particular it is never assumed concave, continuous, or generated by any particular
pension schedule.
